\BeleskeID{04-s027}
% element: 04-s027:e01 — Digresija i osmobitni registar 01101000
% element: 04-s027:e02 — Razvijeni izraz RVr i rezultat skaliranja CVr/r^k
% element: 04-s027:e03 — Tri položaja, zapisi i stvarne vrednosti CV/2^5, CV, CV/2^8; drugi položaji
% element: 04-s027:e04 — Međukorak svođenja razlomljenih težina na zajednički imenilac
% element: 04-s027:d01

Svaku težinu sa negativnim eksponentom možemo prepisati sa zajedničkim imeniocem:
\[\begin{aligned}
RV_r&=c_{-1}r^{-1}+c_{-2}r^{-2}+\cdots+c_{-k+1}r^{-k+1}+c_{-k}r^{-k}\\
&=\frac{c_{-1}r^{k-1}+c_{-2}r^{k-2}+\cdots+c_{-k+1}r^1+c_{-k}r^0}{r^k}\\
&=\frac{CV_r}{r^k}.
\end{aligned}\]
Brojilac je vrednost istog niza cifara kada ga čitamo kao celobrojni zapis. Isti princip skaliranja primenjujemo na celu reč registra: ako je $k$ bitova desno od tačke, neoznačenu celobrojnu vrednost delimo sa $2^k$.

Za sadržaj $01101000$ celobrojna vrednost je $104$. Tačka iza $D_5$ daje $104/2^5=3.25$; iza $D_0$ daje $104$; ispred $D_7$ daje $104/2^8=0.40625$. Različite vrednosti nastaju zbog položaja tačke, bez promene bitova.
